% ============================================================
% SUPPLEMENT 5: Goodfire / Neural Geometry Integration
% The Geometry of SAE Artifacts Research Program
% May 2026
% ============================================================
\documentclass[11pt,a4paper]{article}

\usepackage[utf8]{inputenc}
\usepackage[T1]{fontenc}
\usepackage{times}
\usepackage{amsmath,amssymb,amsfonts,amsthm}
\usepackage{graphicx}
\usepackage{booktabs}
\usepackage{hyperref}
\usepackage[margin=1in]{geometry}
\usepackage{enumitem}
\usepackage{caption}
\usepackage{subcaption}
\usepackage{natbib}
\usepackage{xcolor}
\usepackage{microtype}
\usepackage{float}
\usepackage{fancyhdr}
\usepackage{multirow}
\usepackage{array}
\usepackage{longtable}
\usepackage{tcolorbox}
\usepackage{colortbl}

\graphicspath{{../}{../figures/}}
\hypersetup{colorlinks=true,linkcolor=blue!60!black,citecolor=blue!60!black,urlcolor=blue!60!black}

\newtcolorbox{keyinsight}{colback=blue!5!white,colframe=blue!50!black,title=Idea clave,fonttitle=\bfseries}
\newtcolorbox{finding}[1]{colback=green!5!white,colframe=green!50!black,title=#1,fonttitle=\bfseries}
\newtcolorbox{warningbox}{colback=orange!5!white,colframe=orange!70!black,title=Advertencia,fonttitle=\bfseries}
\newtcolorbox{hypobox}[1]{colback=purple!5!white,colframe=purple!50!black,title=#1,fonttitle=\bfseries}
\newtcolorbox{protocolbox}[1]{colback=gray!5!white,colframe=gray!50!black,title=#1,fonttitle=\bfseries}

\newtheorem{definition}{Definicion}
\newtheorem{hypothesis}{Hypothesis}
\newtheorem{proposition}{Proposicion}

\setlength{\parskip}{0.4em}
\setlength{\parindent}{1.2em}

\pagestyle{fancy}
\fancyhf{}
\rhead{\small\textit{Supplement 5: Goodfire / Neural Geometry}}
\lhead{\small\textit{Geometry of SAE Artifacts --- May 2026}}
\cfoot{\thepage}

\newcommand{\cosED}{\cos(E_i,D_i)}
\newcommand{\Ef}{E_f}
\newcommand{\Df}{D_f}
\newcommand{\Mh}{\mathcal{M}_h}
\newcommand{\My}{\mathcal{M}_y}
\newcommand{\Acal}{\mathcal{A}}
\newcommand{\Bcal}{\mathcal{B}}
\newcommand{\Hellinger}{d_{\mathrm{H}}}
\DeclareMathOperator*{\argmax}{arg\,max}

\title{\Large\bfseries Supplement 5 Report: Goodfire Neural Geometry Integration\\[0.4em]
\large The Geometry of SAE Artifacts Research Program\\[0.3em]
\normalsize From Sparse Linear Features to Representation Manifolds, SAE Shattering, and Manifold-Aware Control}
\author{\textit{Working Report --- May 2026}\\[0.3em]
\small Integrating Goodfire's \emph{The World Inside Neural Networks} and recent manifold literature}
\date{}

\begin{document}
\maketitle
\thispagestyle{fancy}

\begin{abstract}
\noindent This supplement integrates the Goodfire Neural Geometry program, especially \emph{The World Inside Neural Networks}~\citep{geiger2026worldinside} and the associated manifold-steering paper~\citep{wurgaft2026manifoldsteering}, into the Geometry of SAE Artifacts research program. The central update is conceptual: Sparse Autoencoder (SAE) features should not be treated as the final ontology of a model. Many model-internal concepts are better described as structured geometric objects---curves, circles, helices, regions, subspaces, hierarchies, or manifolds---whose local linear pieces may be fragmented by an SAE dictionary into many apparently unrelated features.

We introduce a new hypothesis, \textbf{H7: Manifold Fragmentation / SAE Shattering}. H7 does not replace the already-supported geometric hedge mechanism from Phase~2 and Supplement~3; instead it explains how local anti-aligned features can arise as corrective shards of a higher-dimensional representation manifold. This reframes anti-alignment, the F1803 temporal manifold, and Neuronpedia logit inversion as instances of a broader neural-geometry problem: the mismatch between curved representation geometry and global sparse linear dictionaries.

This document is explicitly non-retroactive. It does not reopen closed gate decisions G0--G2 or revise the empirical claims of Supplements~3--4. It provides a theory bridge, a literature update, and a concrete experimental protocol for Supplement~5 and later Phase~3--5 work.
\end{abstract}

\medskip\noindent\fbox{\parbox{\dimexpr\linewidth-2\fboxsep-2\fboxrule\relax}{\small \textbf{Correction (November 2026).} This report cited results from the March~2026 version of Supplement~4, which has been withdrawn and replaced. In particular, the ``active anti-aligned features'' it mentioned were randomly selected layer-8 features with ordinary encoder--decoder alignment, and the revised Supplement~4 finds that every strongly anti-aligned GPT-2 feature queried is dead or nearly dead. The affected statements below have been revised. See \texttt{docs/ERRATA.md} (E6, E7) in the repository.}}

\vspace{0.5em}\hrule\vspace{0.5em}
\tableofcontents
\newpage

\section{Non-Retroactivity and Source Corpus}
\label{sec:nonretro}

\begin{warningbox}
This supplement updates the \emph{interpretive frame} of the project, not the historical record of completed phases. PHASE0, PHASE1, PHASE2, SUPPLEMENT3, and SUPPLEMENT4 remain closed reports. Their results should be cited as measured. New neural-geometry claims belong in this supplement and in future experiments.
\end{warningbox}

The source corpus consisted of locally extracted text from the following Goodfire pages (not redistributed with this repository):
\begin{itemize}[noitemsep]
  \item \texttt{goodfire\_world\_inside\_neural\_networks\_extracted.txt}
  \item \texttt{manifold\_steering\_extracted.txt}
  \item \texttt{neural\_geometry\_series\_extracted.txt}
  \item \texttt{phylogeny\_manifold\_extracted.txt}
  \item \texttt{alzheimers\_detection\_extracted.txt}
\end{itemize}

The primary link requested for integration is:
\begin{center}
\url{https://www.goodfire.ai/research/the-world-inside-neural-networks}
\end{center}
It was published as part of Goodfire's Neural Geometry Series in May~2026.

\section{The Goodfire Neural Geometry Framework}
\label{sec:goodfire}

\subsection{Core Claim}

Goodfire's central claim is that neural networks learn structured inner worlds. These inner worlds are not merely collections of isolated linear directions; they contain geometric structures that recapitulate the structure of the external world represented in the training data. Examples include cyclic calendar structures, number helices, smooth historical-year curves, color surfaces, spatial layouts, phylogenetic trees, and biomedical manifolds~\citep{geiger2026worldinside,pearce2025treeoflife,wang2026alzheimers}.

\begin{keyinsight}
The key shift is from \emph{feature ontology} to \emph{geometric ontology}. A concept may be a point, a direction, a curve, a circle, a helix, a region, a graph, a hierarchy, or a manifold. A single SAE latent is one possible local coordinate, not necessarily the concept itself.
\end{keyinsight}

Goodfire frames interpretability across three coupled levels:
\begin{enumerate}[noitemsep]
  \item \textbf{Representations}: geometric data structures in activation space.
  \item \textbf{Computations}: transformations that consume and produce those structures.
  \item \textbf{Behavior}: output distributions or actions generated from those computations.
\end{enumerate}
In this framing, representation geometry is analogous to a model's internal data structure. Understanding a computation without understanding the representation it operates on is incomplete.

\subsection{Representation and Behavior Manifolds}

Let $\Acal_\ell \subset \mathbb{R}^{d}$ denote the activation space at layer $\ell$. A concept $c$ may occupy a low-dimensional submanifold
\[
  \Mh(c) \subset \Acal_\ell.
\]
For language models, the behavior space can be treated as a probability simplex over vocabulary distributions:
\[
  \Bcal \subset \Delta^{|V|-1}.
\]
Goodfire's manifold-steering work defines a corresponding behavior manifold $\My(c) \subset \Bcal$ and studies whether causal interventions along $\Mh(c)$ induce natural trajectories along $\My(c)$~\citep{wurgaft2026manifoldsteering}.

For output distributions $p$ and $q$, the behavior geometry can be measured with the Hellinger distance:
\[
  \Hellinger(p,q) = \frac{1}{\sqrt{2}}\left\|\sqrt{p}-\sqrt{q}\right\|_2.
\]

\subsection{Linear Steering as a Local Approximation}

Conventional activation steering uses a vector
\[
  h'_\ell = h_\ell + \alpha v,
\]
where $v$ is often a difference of means between positive and negative examples. Goodfire argues that this works when the concept geometry is locally flat or globally linear. If the concept is curved, a straight chord can leave the natural activation manifold and produce incoherent behavior.

Manifold steering replaces the straight-line intervention with a path
\[
  \gamma(t) \in \Mh(c),
\]
and evaluates whether
\[
  F_\ell(\gamma(t)) \approx \eta(t), \qquad \eta(t) \in \My(c),
\]
where $F_\ell(h)=p_\theta(\cdot \mid do(h_\ell=h))$ is the causal activation-to-behavior map.

\section{Why This Matters for SAEs}
\label{sec:sae}

\subsection{The SAE Shattering Problem}

Goodfire explicitly notes that SAEs can fragment a smooth manifold into many local, apparently unrelated features. In their rhyming example, an unsupervised manifold captures phonological endings as a smooth structure; SAE auto-interpretability labels instead identify local fragments that miss the global concept~\citep{geiger2026worldinside}.

\begin{definition}[SAE manifold shattering]
An SAE shatters a concept manifold $\Mh(c)$ when the global concept is represented not by a single latent, but by a collection of local latents $\{f_1,\ldots,f_k\}$ whose activation regions cover patches of $\Mh(c)$ while obscuring the manifold's topology or intrinsic coordinates.
\end{definition}

This is not a claim that SAEs are useless. Rather, SAEs are useful probes with architectural priors. The danger is treating the SAE dictionary as the model's true ontology.

\subsection{Architectural Priors Are Geometric Priors}

Recent work makes this point precise. \citet{hindupur2025projecting} argue that each SAE architecture projects assumptions about concept geometry: linear separability, sparsity, homogeneity of intrinsic dimension, and global dictionary structure. \citet{menon2024formal} show in formal-language settings that SAE latents can be interpretable but not causally faithful. \citet{muchane2025hierarchical} introduce hierarchical SAE variants to represent semantic relationships explicitly.

Thus, the current project should distinguish:
\begin{itemize}[noitemsep]
  \item \textbf{Model geometry}: the actual representation structure learned by the transformer.
  \item \textbf{SAE geometry}: the geometry imposed by the dictionary learning objective.
  \item \textbf{Analysis geometry}: the geometry visible through PCA, UMAP, kNN graphs, local factor models, or manifold fitting.
\end{itemize}

\section{Updated Literature Map}
\label{sec:lit}

\begin{longtable}{@{}p{3.2cm}p{5.2cm}p{6.2cm}@{}}
\caption{Recent literature relevant to the Goodfire update.}\\
\toprule
\textbf{Work} & \textbf{Contribution} & \textbf{Impact on this project}\\
\midrule
\endfirsthead
\toprule
\textbf{Work} & \textbf{Contribution} & \textbf{Impact on this project}\\
\midrule
\endhead
\citet{engels2024multidimensional} & Multidimensional language-model features; days and months as circles; causal tests of circular features. & F1803 and temporal features should be treated as possible manifolds, not just directions.\\
\citet{modell2025origins} & Formalizes representation manifolds as metric/topological objects and studies their origins. & Provides language for intrinsic dimension, topology, geodesic distance, and concept metrics.\\
\citet{karkada2026symmetry} & Shows that symmetries in language statistics shape representation geometry. & Suggests temporal manifolds arise from co-occurrence symmetries and latent world structure.\\
\citet{prieto2026correlations} & Shows that correlations shape superposition; interference can be constructive. & Reframes SAE artifacts as interaction between data correlations and compression geometry.\\
\citet{wurgaft2026manifoldsteering} & Demonstrates that steering along representation manifolds yields natural behavior trajectories. & Provides an experimental template for F1803 and anti-aligned-feature steering.\\
\citet{shafran2026regions} & Replaces global directions with local regions using Mixture of Factor Analyzers. & Strong baseline for local-geometry analysis of anti-aligned neighborhoods.\\
\citet{hindupur2025projecting} & SAE architectures impose assumptions about concept geometry. & Supports the claim that anti-alignment is architecture-conditioned projection geometry.\\
\citet{menon2024formal} & SAE latents can be interpretable but not causally faithful in formal languages. & Reinforces causal validation as a requirement before mechanistic claims.\\
\citet{muchane2025hierarchical} & Hierarchical SAEs encode semantic hierarchy explicitly. & Distinguishes pathological shattering from useful hierarchical decomposition.\\
\citet{kantamneni2025trigonometry} & Addition uses helical/trigonometric number representations. & Shows manifolds can be computational substrates, not just visualizations.\\
\citet{levy2024digits} & Numbers represented via base-10 circular digit structures. & Suggests temporal/numeric features may be product manifolds or periodic coordinates.\\
\citet{fu2026convergent} & Distinguishes spectral similarity from functional geometric separability. & Warns that visible geometry must be tested causally before being called a mechanism.\\
\bottomrule
\end{longtable}

\section{Integration With Existing Results}
\label{sec:integration}

\subsection{Phase 0--2 and Supplements 3--4}

The completed project phases established the following:
\begin{itemize}[noitemsep]
  \item \textbf{Phase 0--1}: anti-aligned features exist cross-model and are distributionally structured.
  \item \textbf{Phase 2}: anti-alignment appears even with $\lambda=0$ and under TopK/JumpReLU variants, supporting a geometric origin rather than a pure L1 artifact.
  \item \textbf{Supplement 3}: real GPT-2 L8 features pass a bisector/hedge test; anti-alignment has a phase transition around $5$--$6\times$ expansion, and at $32\times$ the toy model matches the real anti-aligned count.
  \item \textbf{Supplement 4 (revised)}: hedge and parent decoders are nearly antiparallel, so their logit effects are inverted; all strongly anti-aligned GPT-2 features queried, and their parents, are dead or nearly dead. The only active anti-aligned feature found is Gemma~2B F3378.
\end{itemize}

Goodfire does not invalidate these results. It adds an interpretive layer: the ``hedge'' may be a local linear correction created when an SAE dictionary approximates a curved or multidimensional representation with sparse linear shards.

\subsection{F1803 as the Most Direct Bridge}

The F1803 temporal-specificity project is the strongest bridge between the current project and neural geometry. Prior measurements found:
\begin{itemize}[noitemsep]
  \item F1803 acts as a temporal specificity gate for post-2019 year contexts.
  \item The temporal representation is not a simple 1D line; PCA over temporal prompts gave a high $\mathrm{PC2}/\mathrm{PC1}$ ratio ($0.87$), consistent with a multidimensional manifold.
  \item The 2023 anomaly was explained by GPT-2 BPE tokenization: `` 2023'' decomposes as `` 20'' + ``23'', while `` 2024'' is a dedicated token.
  \item L0\_MLP behaves as a multichannel router rather than a universal hub.
\end{itemize}

Under the Goodfire framing, F1803 should be analyzed as a local coordinate or shard of a temporal manifold. The key question becomes: what is the intrinsic topology and metric of the temporal representation, and how does F1803 sit relative to its tangent and normal directions?

\subsection{Neuronpedia Logit Inversion as Local Geometry}

Supplement~4 found strong logit inversion between hedge and parent features, which follows from their nearly antiparallel decoders. In the old framing, this geometry is read as subtractive hedging; note that the GPT-2 hedges and parents examined are dead on Neuronpedia's activation sample, so the corrective role is a description of the weights, not of observed inference-time behaviour. In the neural-geometry framing, it can be interpreted more generally as a local corrective chart: the anti-aligned feature detects entry into a region of representation space and writes a vector that keeps the reconstruction on or near the learned manifold.

\begin{keyinsight}
Anti-alignment may be a \emph{local normal correction}: the encoder detects a region or parent cluster; the decoder writes a displacement that cancels a local off-manifold or over-represented component. This is compatible with the hedge-family mechanism but embeds it in a richer manifold view.
\end{keyinsight}

\section{New Hypothesis H7}
\label{sec:h7}

\begin{hypobox}{H7: Manifold Fragmentation / SAE Shattering}
Anti-aligned SAE features are sometimes local linear shards of a curved or multidimensional concept manifold. The SAE objective, especially at large expansion ratios, decomposes the manifold into sparse dictionary elements. Some shards become corrective: their encoder detects a local region or parent cluster, while their decoder writes an opposing vector that keeps reconstruction consistent with the global manifold.
\end{hypobox}

H7 predicts:
\begin{enumerate}[noitemsep]
  \item Anti-aligned features should lie in dense decoder-neighborhoods or coactivation neighborhoods rather than being isolated random outliers.
  \item For a concept with known topology, anti-aligned features should occupy high-curvature or transition regions of the reconstructed manifold.
  \item The chord distance between feature shards should underestimate geodesic distance when the concept is curved.
  \item Local PCA or factor analysis around anti-aligned neighborhoods should reveal dimension greater than one.
  \item Linear steering through anti-aligned-feature directions should be less natural than manifold-aware steering along local tangent directions.
  \item Architectures with explicit hierarchy or local regions should reduce apparent shattering while preserving the concept's global geometry.
\end{enumerate}

H7 differs from existing hypotheses as follows:
\begin{itemize}[noitemsep]
  \item H1 (Subtractive Hedges) is a local mechanism; H7 is a global geometric setting in which H1 can occur.
  \item H2 (Dead Feature Decay) predicts low activation and little structure; H7 predicts structured neighborhoods even when some shards are dead.
  \item H3/H5 (Superposition/Compositionality) predict distributed structure; H7 specifies that the distributed structure may be a manifold with intrinsic topology.
  \item H4 (Ghost Gradients) is an optimization-dynamics explanation; H7 can operate without ghost gradients.
\end{itemize}

\section{Supplement 5 Experimental Protocol}
\label{sec:protocol}

\subsection{S5A: Geodesic-vs-Chord Test}

Construct kNN graphs over decoder directions, encoder directions, and activation-pattern similarities for a target family (e.g. F19139/F18573/F11969 plus parents; F1803 temporal cluster). Compare:
\[
  \rho_{\mathrm{geo/chord}}(i,j)=\frac{d_{\mathrm{graph}}(i,j)}{\|D_i-D_j\|_2}.
\]
Large ratios indicate curved or non-convex geometry where straight-line interpolation is misleading.

\subsection{S5B: Local Tangent/Normal Decomposition}

For a target feature $f$, collect nearest neighbors $N_k(f)$ by decoder cosine, encoder-to-decoder similarity, and coactivation. Fit local PCA to their activations or decoder vectors. Decompose $E_f$ and $D_f$ into tangent and normal components:
\[
  E_f = E_f^{\parallel}+E_f^{\perp}, \qquad D_f = D_f^{\parallel}+D_f^{\perp}.
\]
H7 predicts anti-alignment may concentrate in the normal component even when tangent components remain semantically coherent.

\subsection{S5C: SAE Shard Reconstruction}

Given a conceptual trajectory (years, months, weekdays, rhymes, or temporal prompts), identify active SAE features and test whether the active-feature graph forms a path, cycle, surface, or branching structure. Compute:
\begin{itemize}[noitemsep]
  \item intrinsic dimension via local PCA / participation ratio;
  \item continuity of active feature sets along the conceptual variable;
  \item whether auto-interp labels are local while the global graph is semantically coherent.
\end{itemize}

\subsection{S5D: Manifold Steering vs Linear Steering}

For F1803 or a controlled calendar concept, compare:
\begin{enumerate}[noitemsep]
  \item linear steering along a single decoder or contrast vector;
  \item tangent steering along a fitted local manifold;
  \item geodesic steering along a kNN graph or spline fit.
\end{enumerate}
Metrics should include target-token probability, off-target probability, Hellinger distance to natural behavior trajectories, and distance to natural activations.

\subsection{S5E: Alternative Decompositions}

Compare SAEs against local-region methods such as Mixture of Factor Analyzers~\citep{shafran2026regions}, hierarchical SAEs~\citep{muchane2025hierarchical}, or geometry-aware variants such as SpaDE-style architectures~\citep{hindupur2025projecting}. The goal is not to replace SAEs, but to diagnose what SAEs miss.

\section{Implications for Future Phases}
\label{sec:future}

\subsection{Phase 3: Seed Stability}

Seed stability should track not only exact feature identity, but also neighborhood identity. If a specific anti-aligned feature changes across seeds while the same local manifold region and parent cluster persist, the correct conclusion is not ``unstable artifact'' but ``stable geometric region with seed-dependent discretization.''

\subsection{Phase 4: Coactivation Case Studies}

Coactivation should be analyzed as a graph or manifold, not merely as pairwise parent-child PMI. H7 predicts families of shards with smooth transitions, local neighborhoods, and possible tangent/normal decomposition.

\subsection{Phase 5: Architectural Comparison}

Architectural comparison should ask: which SAE architectures fragment manifolds, which preserve hierarchy, and which recover local charts? Expansion ratio, threshold type, and sparsity mechanism should be interpreted as geometric priors.

\section{Limitations}
\label{sec:limitations}

\begin{warningbox}
Goodfire's framing is highly relevant and current, but not direct evidence about the Bloom GPT-2 SAE unless we run the proposed tests. This supplement therefore records a theory update and experimental agenda, not a completed empirical confirmation of H7.
\end{warningbox}

Additional caveats:
\begin{itemize}[noitemsep]
  \item Many 2026 references are recent arXiv preprints and should be treated as developing literature.
  \item Low-dimensional plots can mislead; every manifold claim should be paired with causal or metric validation.
  \item Manifold steering provides local control evidence, not global alignment guarantees.
  \item SAE shattering can be pathological or useful: hierarchical decomposition may intentionally break a concept into meaningful levels.
  \item All strongly anti-aligned GPT-2 features examined so far are dead or near-dead (Supplement~4); H7 must account for this inactive population, and the single active example found (Gemma~2B F3378) before generalising.
\end{itemize}

\section{Conclusion}
\label{sec:conclusion}

The Goodfire Neural Geometry framework strengthens the central direction of this research program. Phase~2 and Supplement~3 already showed that anti-alignment is geometric rather than a pure L1 artifact. Goodfire and the 2024--2026 manifold literature show why this is not an isolated oddity: neural representations often organize as curved, structured objects, and sparse linear dictionaries can fragment those objects into local shards.

The immediate conceptual update is therefore:
\begin{quote}
\emph{SAE artifacts are not merely failures of dictionary learning; they are diagnostic traces of the mismatch between sparse linear decompositions and the richer geometry of model representations.}
\end{quote}

The next empirical step is to test H7 with geodesic, local-PCA, shard-reconstruction, and manifold-steering experiments. If confirmed, this will connect the project's anti-alignment results to a broader theory of how neural networks represent and compute over the world inside their activations.

\newpage
\bibliographystyle{plainnat}
\bibliography{goodfire_neural_geometry}

\end{document}
